%%=============================================================================
%% Samenvatting nederlands
%%=============================================================================

\chapter*{Samenvatting nederlands}


In mijn tweede en laatste jaar van de opleiding Programmeren aan HOGENT 
moest ik een graduaatsproef maken. Dit moest een project zijn in samenwerking 
met mijn stagebedrijf. Samen met Victor Buyck Steel Construction, een 
staalbouwbedrijf in Eeklo, hebben we gekozen om een onthaalapplicatie te 
ontwikkelen.

Het bedrijf betaalde jaarlijks voor een externe SaaS-oplossing voor hun 
onthaal. Het probleem was dat deze oplossing regelmatig uitviel en dat de 
jaarlijkse kost hoog opliep. Daarom werd beslist om dit zelf te ontwikkelen.

Het resultaat is WelcomeDesk, een onthaalapplicatie bestaande uit een 
Blazor-frontend en een C\#-backend. Bezoekers kunnen zich aanmelden door 
hun gegevens in te vullen en een medewerker te selecteren. Die medewerker 
wordt dan automatisch verwittigd via een Teams-bericht, een sms en een 
e-mail.

Voor de ontwikkeling werd gekozen voor Blazor omdat het bedrijf hier al 
ervaring mee had en het actief gebruikt in andere projecten. Een 
bijkomend voordeel van Blazor WebAssembly is dat de applicatie als een 
standalone kan draaien zonder dat er een constante 
serververbinding nodig is. Dit maakt de applicatie stabieler en makkelijker 
te deployen op een toestel in het onthaal. De gegevens worden opgeslagen 
in een MySQL-database. De Teams-notificaties verlopen via Power Automate, 
de sms via Proximus.

WelcomeDesk is een werkende vervanger voor de oude betaalde oplossing. 
De licentiekost valt weg, het systeem is stabieler en het kan in de toekomst 
makkelijk uitgebreid worden indien nodig.
